\part{Some Language and Notation of Mathematics}

\section{Sets}

\subsection{Basic Set Theory}

\begin{defbox}
   A \textbf{set} is a collection of elements.
\end{defbox}

Intuitively, a set can be thought of as a bag---our set---which may or may not contain stuff---our elements. Sets are the most fundamental objects in mathematics and many important concepts, such as functions and spaces, can be defined using sets.

If we want to say that an element $x$ belongs to some set $S$, we write
\[x \in S\]
with the ``$\in$" symbol meaning ``an element of" (so $x \in S$ literally means $x$ is an element of $S$). 

When defining a set, we use curly braces ``$\{\}$" to enclose the elements inside of it.

\begin{example}
   The following are some basic examples of sets:
\begin{enumerate}
   \item $A := \{0, 3, \heartsuit\}$ defines a set, named $A$, to the elements $0$, $3$, and $\heartsuit$ (elements do not necessarily need to be numbers, indeed another set can also be an element!) 
   \item $B := \{\text{all odd integers}\}$ defines a set, named $B$, to all odd integers (e.g. $-1 \in B$ and $127 \in B$).
   \item $C := \{n \colon n \text{ is an integer satisfying } n^3 = n\}$ (the colon ``:" in between the curly braces stands for ``such that," so $C$ is the set of all integers $n$ such that $n^3 = n$ (equivalently, $C := \{-1, 0, 1\}$))
   \item The standard sets:
     \begin{flalign*}
         & \NN := \{0\footnotemark[1], 1, 2, 3, \dots\} \text{  (the set of natural numbers)} & \\
         & \ZZ := \{\dots, -2, -1, 0, 1, 2, \dots\} \text{  (the set of integers)} & \\
         & \QQ := \left\{\displaystyle\frac{a}{b} \colon a, b \in \ZZ, b \neq 0\right\} \text{  (the set of rational numbers)} & \\
         & \RR := \text{the set of real numbers}\footnotemark[2] & \\
         & \CC := \left\{a + bi \colon a, b \in \RR, i = \sqrt{-1}\right\} \text{  (the set of complex numbers)} 
      \end{flalign*} 
\end{enumerate}
\footnotetext[1]{Some sources omit $0$ from the set of natural numbers, and define $\NN := \{1, 2, 3, \dots\}$. However in this class, we will always assume $0 \in \NN$, and define $\NN^+ := \{1, 2, 3, \dots\}$.}
\footnotetext[2]{The set of real numbers is surprisingly difficult to define. Take MATH255 to learn how to rigorously define $\RR$.}
\end{example}

\begin{remark}
   Order and repetition in sets do not matter. Indeed, the set $A := \{0, 3, \heartsuit\}$ is equal to the sets $A_1 := \{3, \heartsuit, 0\}$ and $A_2 := \{0, 3, \heartsuit, 0\}$.
\end{remark}

\begin{defbox}
   Two sets $A$ and $B$ are said to be equal if and only if they have the same elements.

   Symbolically, we write $A = B \iff \forall x, x \in A \text{ when } x \in B$.
\end{defbox}

In the above theorem, we introduced two key symbols often used in mathematics. The first, the ``$\iff$" symbol, stands for ``if and only if." In mathematics, this means that the statement on the left is equivalent to the statement on the right. The second symbol, the  ``$\forall$", stands for ``for all" or ``for every" and is called the \textbf{universal quantifier}. Symbolic notation is often eschewed in formal writing but is used very often in more casual mathematical writing.

\vspace{0.2cm}

\begin{remark}
  By Definition 1.0.3, we see that the set $A$ in Remark 1.0.2 is equal to the sets $A_1$ and $A_2$. 
\end{remark}

\begin{defbox}
   Let $A$ and $B$ be arbitrary sets. We say that $A$ \textbf{is a subset of} $B$, or $A$ \textbf{is contained in} $B$, if for every $x \in A$, $x \in B$. In this case we write, $A \subseteq B$, where ``$\subseteq$" is the subset symbol\footnotemark[3].
\end{defbox}

\footnotetext[3]{Many textbooks use the symbol ``$\subset$" to denote a subset instead of $\subseteq$, i.e. if $A$ is a subset of $B$, then $A \subset B$ is written. In these notes, I will use the symbol $\subset$ only when one set is a \textit{proper subset} of the other, in the same way that $<$ defines a strict inequality.}


\vspace{0.2cm}

\begin{remark} 
   Note that $A = B$ holds when $A \subseteq B$ and $B \subseteq A$ both hold. This gives us a very useful way to prove if two sets are equal (we simply prove that one set is the subset of the other and vice versa). 
\end{remark}

\begin{defbox}
   The set with no elements is called the \textbf{empty set}, and is denoted $\emptyset := \{\}$.
\end{defbox}

Intuitively, the empty set can be thought of as a bag with nothing inside of it. The empty set is \textit{not} nothing, but it contains nothing.

\begin{example}
   We conclude with some examples of sets and some of their properties:
   \begin{enumerate}
      \item For every set $A$, $\emptyset \subseteq A$.
   \item $\emptyset \neq \{\emptyset\}$, indeed we see that $\emptyset \subseteq \{\emptyset\}$ but $\{\emptyset\} \not\subseteq \emptyset$. 
   \item $\{\emptyset, \{\emptyset\}\} \neq \{\{\emptyset\}\}$
   \item In set theory, we write the numbers $0 := \emptyset$, $1 := \{0\}$, $2 := \{0,1\}$, and so on.
   \end{enumerate}
\end{example}


